\section{The collision between two efficient pairs}\label{app:full-collision}

This appendix proves Lemma~\ref{lem:two-pairs}. Retain the parameters
$n=3m+2$, $p_0=(m+1)(2m+1)$, $v=2m(m+1)$, $b=p_0-1$, and $c=v-2$.

\begin{proof}
Normalize the first short-run entry to $(3,\ldots,n,2,1,0)$.
In a canonical $C$ cell, the nonzero cyclic core of its $r$th full run is
\[
 L_r=[\operatorname{rot}^{r-1}(1,3,\ldots,n-1,2),n],
 \qquad 1\le r\le n-3.
\]
Representatives beginning in $1$ give
\[
 L_1=[1,3,\ldots,n-1,2,n],\qquad
 L_{n-3}=[1,3,\ldots,n-3,n,n-2,n-1,2].
\]
Set $M=h_ch_b$ and $j=v-m-4$. Number cells from zero. The last full
run of cell $j$ in the second chain of the first pair has actual core
$h_b^{-1}g^{-j}L_{n-3}$. The first full run of cell $m$ in the second
chain of the second pair has core $M^{-1}h_b^{-1}g^{-m}L_1$.
Both cells precede their chains' final $B$ cell, since
$0\le j\le c-2$ and $m\le c-2$.

Equality of these two cores is equivalent to $EL_{n-3}=L_1$, where
\[
 E=g^mh_bh_cg^{-j}=g^m\rho g^{-1}\rho g^{m+2}.
\]
The last expression uses $g^{p_0}=1$, $c=p_0-m-3$ and
$j=p_0-2m-5$. The two nonzero cycles of $g$ are
\[
 (1,5,8,\ldots,3m+2),\qquad
 (2,4,7,\ldots,3m+1,3,6,\ldots,3m).
\]
Applying the displayed expression for $E$ gives
\[
\begin{gathered}
 E(1)=3,\quad E(2)=1,\quad E(x)=x+1\ (3\le x\le n-3),\\
 E(n-2)=2,\quad E(n-1)=n,\quad E(n)=n-1.
\end{gathered}
\]
It sends $L_{n-3}$ to $[3,4,\ldots,n-1,2,n,1]$, a rotation of $L_1$.
Thus the two full runs visit the same two-loop, contradicting class
distinctness.
\end{proof}

\subsection{The initial case on nine symbols}

At $m=2$, the efficient tiles are $(12)$, $(13,11)$ and $(14,10)$.
Their transfers lie in a group of order $12$. To see this directly, put
$A=(2\ 8)$, $Q=(2\ 8\ 7)$, $U=(1\ 3)$ and $V=(5\ 6)$.
The three transfers are $AV,AU,QUV$. They belong to
\[
 \{\sigma U^eV^f:
 \sigma\in\mathfrak S_{\{2,7,8\}},\ e,f\in\{0,1\},\
 \operatorname{sgn}(\sigma)=(-1)^{e+f}\},
\]
a group of order $12$. Twelve consecutive efficient interior tiles
would give thirteen prefix maps in this group, two equal, and hence
a repeated start class. Positive losses are integers and the final
allowance is two. Lemma~\ref{lem:tiles} therefore gives
$\CP(15,24-1/12)=\CP(15,287/12)$.

